% Lecture visuelle directe avec la vision du modèle ; aucun OCR.
% Source : PDF original, vue 002, côté droit, page manuscrite 3.
% Statut : transcrit ; relecture indépendante à effectuer.
\chapter*{Aux membres des Congrégations des hommes et des jeunes gens de Lyon}
\addcontentsline{toc}{chapter}{Aux membres des Congrégations des hommes et des jeunes gens de Lyon (8 mai 1841)}
\markboth{Aux membres des Congrégations}{Aux membres des Congrégations}

\begin{center}
\textit{8 mai 1841, veille du 4\up{e} dimanche après Pâques.}
\end{center}

Messieurs, mes amis et mes frères,

Nous allons célébrer demain cette fête solennelle d'actions de grâces, qui fut établie parmi nous en 1804. C'est le jour que nous avons choisi pour remercier chaque année le Seigneur des grâces qu'il ne cesse de répandre sur la Congrégation et sur les Congréganistes. Quoique je sois privé d'assister de corps à cette douce réunion, mon âme n'y sera pas moins présente. J'espère que, comme à l'ordinaire, nos cœurs s'uniront dans la fraction du pain pour se réchauffer ensemble au feu ardent de l'amour divin. Ma bouche prononcera aussi avec vous ce cantique d'actions de grâces destiné par l'Église à célébrer les bienfaits que le Ciel verse constamment sur nous, comme une douce rosée.

C'est le jour de cette fête si touchante pour des cœurs qui ne sont pas ingrats, que je crois devoir choisir pour poser la première pierre d'un monument que je désirais depuis longtemps voir élever, afin qu'il pût rappeler à jamais la mémoire de tant de bienfaits. Je forme le vœu qu'avec la grâce du Seigneur ce monument subsiste, qu'il demeure au milieu de nous
% Source : PDF original, vue 003, côté gauche, page manuscrite 4.
et qu'il fasse connaître à tous ceux qui nous suivront les moyens admirables par lesquels le Dieu de bonté nous a préparé et nous conserve un abri pour éloigner de nous les dangers qui nous environnent. C'est là qu'il daigne soutenir et guider dans la voie du salut nos pas mal assurés. Lorsque à l'avenir de jeunes frères seront admis dans la famille, et qu'à la vue de nos usages et de nos règles ils nous adresseront cette demande : que veut dire tout ceci ? Nous pourrons leur répondre : il fut un temps où le Seigneur, par la force de son bras, a miraculeusement fait échapper nos pères d'une Égypte infidèle et incrédule, véritable terre de perdition. Il a placé Marie au milieu d'eux comme l'arche d'alliance qu'il contractait avec eux. Sous ses auspices il les a fait entrer dans une nouvelle terre promise. C'est là que le peuple de Dieu qu'il a appelé et choisi lui-même s'est accru et multiplié ! C'est ainsi qu'il nous a protégés et réchauffés à l'ombre de ses ailes ; qu'il a voulu que nous fussions son héritage, et qu'il a promis d'être lui, pour nous, lui-même notre grande récompense ! Écoutez donc le récit de toutes ces merveilles et chantez avec nous un cantique d'actions de grâces envers le Tout-Puissant, qui dans ce jour vous appelle à partager les faveurs qu'il nous a réservées ! C'est ainsi que sa miséricorde s'étend d'âge en âge sur tous ceux qui le craignent.

Le travail que j'entreprends sera donc une œuvre de reconnaissance que j'offre en votre nom à Dieu notre père, qui nous a prévenus de grâces si spéciales, et à Marie notre auguste protectrice, qui a daigné adopter la
% Source : PDF original, vue 003, côté droit, page manuscrite 5.
Congrégation comme sa famille. Permettez-moi de l'offrir aussi à la Congrégation elle-même, comme un témoignage de ma vive gratitude, car je lui dois tout. Pendant la durée de ma longue carrière et encore aujourd'hui, elle a été toujours ma consolation et mon appui. Je suis maintenant éloigné d'elle. Je ne puis prendre aucune part à ses œuvres. Il m'est doux de pouvoir encore m'occuper d'elle, et de conserver ainsi, au moins de fait si ce n'est de droit, la qualité de membre en activité.

C'est donc avec une douce satisfaction que je commence aujourd'hui, sous les auspices de Marie, à travailler à ce grand ouvrage. Je le dédie et le consacre à Marie, mère de Dieu, Vierge immaculée, reine de la Congrégation. J'ose espérer qu'elle daignera l'agréer et le bénir.

Priez pour moi, mes frères, afin qu'en racontant ce que le Seigneur a fait pour nous, je ne ressemble pas à un airain sonnant ou à une cymbale retentissante. Que mon cœur soit lui-même pénétré des sentimens que je désire faire passer dans les vôtres. Obtenez-moi la grâce de profiter de tant de faveurs que j'ai reçues de la main miséricordieuse de notre Dieu, de persévérer jusqu'à la fin de mes jours dans la carrière de la pénitence, de mériter de trouver grâce et pardon devant le juge souverain, et d'être enfin assez heureux pour rejoindre un jour dans sa gloire ceux qui nous y ont devancés. Puisse le Dieu tout-puissant, père des miséricordes, qui nous a réunis dans le temps, dans la charité, nous y réunir aussi dans l'éternité. Ainsi soit-il.

\begin{center}
\textit{Cor unum et anima una.}\footnote{Actes des Apôtres 4, 32 : « Un seul cœur et une seule âme. » (Note du transcripteur).}
\end{center}
